\chapter{Технологический раздел}

\section{Выбор и обоснование средств реализации}

Профиль нагрузки проектируемой системы сочетает аналитические запросы
на чтение с агрегацией и транзакционную процедуру массового обновления
статистики.

\subsection{СУБД}

Выбор основной СУБД связан с проектными решениями, принятыми в
конструкторском разделе. Расчёт метрик делегирован уровню данных
и реализуется хранимыми функциями в подразделе~\ref{subsec:metric-algorithms},
поэтому требуется процедурный язык на стороне сервера БД. Логическая
проверка командного счёта из раздела~\ref{sec:integrity-design} выполняется
после пакетной загрузки статистики матча, что требует отложенной проверки
целостности к моменту фиксации транзакции. Серверный слой спроектирован
как асинхронный REST API в подразделе~\ref{subsec:app-architecture}, поэтому
взаимодействие с БД не должно блокировать цикл событий приложения. Массовое
обновление сезонной витрины, описанное в подразделе~\ref{subsec:bulk-update},
выполняется одной операцией вставки с обновлением существующих строк.

SQLite~3 является встраиваемой бессерверной СУБД. Такая модель удобна для
локальных приложений и прототипов, но не соответствует серверному сценарию
с несколькими конкурентными подключениями. Кроме того, SQLite не предоставляет
серверный процедурный язык и не имеет штатного асинхронного драйвера для
выбранной модели приложения.

MySQL~8 поддерживает хранимые процедуры и операцию вставки с обновлением.
При этом в рассматриваемом сценарии недостатком является отсутствие
отложенной проверки ограничений целостности к моменту фиксации транзакции.
Это критично для правила, по которому сумма очков игроков должна совпадать
с итоговым счётом команды только после загрузки полного протокола матча.
Также выбранный серверный стек не получает штатного неблокирующего драйвера,
сопоставимого с asyncpg для PostgreSQL.

PostgreSQL~15 удовлетворяет указанным требованиям: поддерживает процедурный
язык PL/pgSQL, отложенную проверку ограничений, атомарную вставку
с обновлением через \texttt{INSERT ... ON CONFLICT DO UPDATE} и работу через
асинхронный драйвер asyncpg~\cite{postgres,asyncpg}. Итоговое сравнение
кандидатов приведено в таблице~\ref{tab:dbms-compare}.

\begin{table}[H]
\caption{Сравнительная характеристика СУБД}
\label{tab:dbms-compare}
\centering
\small
\renewcommand{\arraystretch}{1.2}
\setlength{\tabcolsep}{5pt}
\begin{tabular}{|>{\RaggedRight\arraybackslash}p{0.48\textwidth}
  |>{\centering\arraybackslash}p{0.14\textwidth}
  |>{\centering\arraybackslash}p{0.14\textwidth}
  |>{\centering\arraybackslash}p{0.14\textwidth}|}
\hline
\textbf{Критерий} & \textbf{PostgreSQL 15} & \textbf{MySQL 8} & \textbf{SQLite 3} \\
\hline
Серверный процедурный язык для хранимой логики & Да & Да & Нет \\
\hline
Отложенная проверка ограничений целостности & Да & Нет & Нет \\
\hline
Штатный асинхронный драйвер для выбранного серверного стека & Да & Нет & Нет \\
\hline
Атомарная вставка с обновлением существующей строки & Да & Да & Да \\
\hline
\end{tabular}
\end{table}

По совокупности требований выбрана PostgreSQL~15: она покрывает как
транзакционные механизмы уровня БД, так и интеграцию с асинхронным
серверным приложением.

\subsection{Кэширование}

Паттерн Cache-Aside реализован на базе in-memory СУБД Redis~7~\cite{redis}.
Выбор обоснован поддержкой операций сканирования и группового удаления по
шаблону, что необходимо для инвалидации связанных ключей при пересчёте
статистики отдельного сезона. Команда \texttt{SETEX} позволяет атомарно
устанавливать значение и время жизни ключа, исключая состояние гонки.

\subsection{Серверный слой}

REST API реализован на Python~3.11 с фреймворком FastAPI~\cite{fastapi}.
Фреймворк построен на стандарте ASGI и модели асинхронного ввода-вывода:
каждый обработчик маршрута является корутиной, что позволяет обслуживать
конкурентные запросы в одном процессе без блокировки потока. Механизм
внедрения зависимостей позволяет декларативно привязать к маршруту сессию
конкретной роли СУБД, что соответствует ролевой модели из
раздела~\ref{sec:roles}: разграничение привилегий происходит на уровне
соединения с PostgreSQL, а не в логике приложения. Для взаимодействия с
СУБД применяется драйвер asyncpg~\cite{asyncpg}, реализующий бинарный
протокол PostgreSQL, в связке с SQLAlchemy~2.0 для конструирования
динамических аналитических запросов.

\subsection{Клиентский слой}

Клиентская часть представляет собой одностраничное веб-приложение
на базе библиотеки React~18~\cite{react} и строго типизированного языка
TypeScript. Выбор данного стека технологий обусловлен компонентным подходом
к разработке пользовательского интерфейса, что позволяет эффективно
переиспользовать элементы отображения: таблицы статистики, графики
и карточки игроков. Использование механизма теневого дерева документа
в React обеспечивает высокую производительность при рендеринге объёмных
наборов данных и динамическом обновлении аналитических дашбордов без
полной перезагрузки страницы.

\section{Реализация структуры базы данных}

\subsection{Создание таблиц и ограничений}

Структура базы данных реализована с помощью DDL-скриптов
(полный листинг приведён в Приложении~А). В листинге~\ref{lst:ddl-gps}
представлен фрагмент создания таблицы фактов \texttt{game\_player\_stats},
иллюстрирующий механизм декларативного обеспечения ссылочной и логической
целостности.

\begin{lstlisting}[language=SQL, label=lst:ddl-gps,
  caption=DDL таблицы \texttt{game\_player\_stats} (фрагмент)]
CREATE TABLE game_player_stats (
    stat_id   SERIAL PRIMARY KEY,
    game_id   INTEGER NOT NULL
                REFERENCES games(game_id) ON DELETE CASCADE,
    player_id INTEGER NOT NULL
                REFERENCES players(player_id) ON DELETE RESTRICT,
    team_id   INTEGER NOT NULL
                REFERENCES teams(team_id) ON DELETE RESTRICT,
    points    SMALLINT NOT NULL DEFAULT 0 CHECK (points >= 0),
    fouls     SMALLINT NOT NULL DEFAULT 0 CHECK (fouls BETWEEN 0 AND 6),
    CONSTRAINT uq_game_player UNIQUE (game_id, player_id)
);
\end{lstlisting}

\subsection{Реализация индексов}

Аналитические индексы (Приложение~А) реализованы на базе структуры B-дерева.
В листинге~\ref{lst:indexes} продемонстрировано создание ключевых индексов
для таблицы фактов.

\begin{lstlisting}[language=SQL, label=lst:indexes,
  caption=Создание индексов аналитических запросов (фрагмент)]
CREATE INDEX idx_gps_player_id
    ON game_player_stats(player_id);
CREATE INDEX idx_gps_team_game
    ON game_player_stats(team_id, game_id);
CREATE INDEX idx_pss_season
    ON player_season_stats(season_id, per DESC NULLS LAST);
\end{lstlisting}

Исследование эффективности данных структур и их влияния на время выполнения
запросов приведено в исследовательском разделе (глава~4).

\subsection{Загрузка тестовых данных}\label{subsec:data-load}

Первоначальное наполнение базы данных выполнено средствами языка Python
с использованием библиотеки \texttt{nba\_api}~\cite{nba-api}, формирующей
HTTP-запросы к публично доступным эндпоинтам сервиса
\texttt{stats.nba.com}~\cite{nba-stats}. Скрипт извлекает архивную
статистику NBA за пять регулярных сезонов 2021/22--2025/26.

Полученные данные нормализуются и загружаются в реляционные таблицы
PostgreSQL пакетами в порядке, соответствующем иерархии внешних ключей:
справочники, затем игроки, матчи, статистика матчей, история команд.
Использование официального источника сохраняет репрезентативное распределение
показателей реальных игроков, что критически важно для корректной верификации
относительных метрик.

По результатам контрольных запросов \texttt{SELECT COUNT(*)} объём ключевых
таблиц после загрузки приведён в таблице~\ref{tab:db-count}. Таблица фактов
\texttt{game\_player\_stats} содержит 159\,652 записей, таблица
\texttt{games}~--- 6\,140 матчей. Это полностью покрывает нефункциональное
требование об объёме данных.

{\footnotesize\setlength{\tabcolsep}{4pt}\renewcommand{\arraystretch}{1.15}
\begin{table}[H]
\caption{Объём загруженных данных (контрольный запрос \texttt{SELECT COUNT(*)})}
\label{tab:db-count}
\centering
\begin{tabular}{|l|r|}
\hline
\textbf{Таблица} & \textbf{Число строк} \\
\hline
\texttt{game\_player\_stats}  & 159\,652 \\
\texttt{games}                 & 6\,140 \\
\texttt{players}               & 5\,126 \\
\texttt{player\_season\_stats} & 2\,897 \\
\texttt{teams}                 & 30 \\
\texttt{positions}             & 5 \\
\texttt{seasons}               & 5 \\
\hline
\end{tabular}
\end{table}}

\subsection{Примеры запросов к базе данных}

Для верификации корректности схемы и демонстрации аналитических
возможностей системы разработан набор SQL-запросов, охватывающий основные
классы операций выборки.

Простой запрос с фильтрацией. Выборка активных игроков позиции
центрового. Текст запроса приведён в Приложении~Б,
листинг~\ref{lst:sql-filter}.
Результат выполнения запроса приведён в таблице~\ref{tab:sql-filter-result}.

{\footnotesize\setlength{\tabcolsep}{4pt}\renewcommand{\arraystretch}{1.15}
\begin{table}[H]
\caption{Результат запроса с фильтрацией по позиции центрового}
\label{tab:sql-filter-result}
\centering
\begin{tabular}{|r|l|l|}
\hline
\textbf{ID} & \textbf{Имя} & \textbf{Фамилия} \\
\hline
23  & Steven  & Adams   \\
25  & Bam     & Adebayo \\
61  & Jarrett & Allen   \\
172 & Deandre & Ayton   \\
260 & Charles & Bassey  \\
361 & Goga    & Bitadze \\
363 & Bismack & Biyombo \\
480 & Tony    & Bradley \\
602 & Thomas  & Bryant  \\
694 & Clint   & Capela  \\
\hline
\end{tabular}
\end{table}}

Многотабличный запрос. Получение результатов матчей с названиями
команд и меткой сезона (текст запроса см. в Приложении~Б,
листинг~\ref{lst:sql-join}).
Результат выполнения запроса (первые 10 строк из 1225) приведён
в таблице~\ref{tab:sql-join-result}.

{\footnotesize\setlength{\tabcolsep}{3pt}\renewcommand{\arraystretch}{1.12}
\begin{table}[H]
\caption{Результат многотабличного запроса по матчам сезона 2025-26}
\label{tab:sql-join-result}
\centering
\begin{tabular}{|l|c|l|r|r|l|c|}
\hline
\textbf{Сезон} & \textbf{Дата} & \textbf{Хозяева} & \textbf{Дом.} & \textbf{Гост.} & \textbf{Гости} & \textbf{OT} \\
\hline
2025-26 & 21.10.2025 & Los Angeles Lakers     & 109 & 119 & Golden State Warriors    & 0 \\
2025-26 & 21.10.2025 & Oklahoma City Thunder  & 125 & 124 & Houston Rockets          & 0 \\
2025-26 & 22.10.2025 & Charlotte Hornets      & 136 & 117 & Brooklyn Nets            & 0 \\
2025-26 & 22.10.2025 & Dallas Mavericks       &  92 & 125 & San Antonio Spurs        & 0 \\
2025-26 & 22.10.2025 & Milwaukee Bucks        & 133 & 120 & Washington Wizards       & 0 \\
2025-26 & 22.10.2025 & New York Knicks        & 119 & 111 & Cleveland Cavaliers      & 0 \\
2025-26 & 22.10.2025 & Boston Celtics         & 116 & 117 & Philadelphia 76ers       & 0 \\
2025-26 & 22.10.2025 & Portland Trail Blazers & 114 & 118 & Minnesota Timberwolves   & 0 \\
2025-26 & 22.10.2025 & Utah Jazz              & 129 & 108 & Los Angeles Clippers     & 0 \\
2025-26 & 22.10.2025 & Orlando Magic          & 125 & 121 & Miami Heat               & 0 \\
\hline
\end{tabular}
\end{table}}

Итоговый запрос с агрегацией и группировкой. Средняя
результативность по позициям за сезон, с фильтрацией игроков
с достаточным игровым временем (текст запроса см. в Приложении~Б,
листинг~\ref{lst:sql-group}).
Результат выполнения запроса приведён в таблице~\ref{tab:sql-group-result}.

{\footnotesize\setlength{\tabcolsep}{4pt}\renewcommand{\arraystretch}{1.15}
\begin{table}[H]
\caption{Средняя результативность по позициям (сезон 2025-26)}
\label{tab:sql-group-result}
\centering
\begin{tabular}{|c|r|r|r|r|r|}
\hline
\textbf{Поз.} & \textbf{Игроков} & \textbf{PTS} & \textbf{REB} & \textbf{AST} & \textbf{PER} \\
\hline
C  & 51  & 10.99 & 7.15 & 2.07 & 20.34 \\
PF & 44  & 13.89 & 5.95 & 2.41 & 18.46 \\
PG & 74  & 13.34 & 2.96 & 4.44 & 16.07 \\
SF & 60  & 11.43 & 4.29 & 2.19 & 14.59 \\
SG & 110 & 11.27 & 3.67 & 2.43 & 14.07 \\
\hline
\end{tabular}
\end{table}}

Запрос с подзапросом. Игроки, чей PER превышает среднее
по лиге за сезон (текст запроса см. в Приложении~Б,
листинг~\ref{lst:sql-subquery}).
Результат выполнения запроса (10 строк из 263) приведён
в таблице~\ref{tab:sql-subquery-result}.

{\footnotesize\setlength{\tabcolsep}{4pt}\renewcommand{\arraystretch}{1.15}
\begin{table}[H]
\caption{Игроки с PER выше среднего по лиге (сезон 2025-26)}
\label{tab:sql-subquery-result}
\centering
\begin{tabular}{|l|c|r|r|}
\hline
\textbf{Игрок} & \textbf{Команда} & \textbf{PER} & \textbf{PTS} \\
\hline
Nikola Jokić              & DEN & 37.57 & 27.68 \\
Giannis Antetokounmpo     & MIL & 36.07 & 27.58 \\
Victor Wembanyama         & SAS & 34.58 & 25.05 \\
Shai Gilgeous-Alexander   & OKC & 31.27 & 31.16 \\
Luka Dončić               & LAL & 30.37 & 32.97 \\
Jalen Duren               & DET & 28.65 & 18.25 \\
Kawhi Leonard             & LAC & 28.24 & 27.89 \\
Joel Embiid               & PHI & 26.88 & 26.18 \\
Ty Jerome                 & MEM & 26.52 & 11.80 \\
Jalen Johnson             & ATL & 26.33 & 21.61 \\
\hline
\end{tabular}
\end{table}}

Составной запрос (UNION). Сводный список лидеров по очкам
и передачам в одной выборке (текст запроса см. в Приложении~Б,
листинг~\ref{lst:sql-union}).
Результат выполнения запроса приведён в таблице~\ref{tab:sql-union-result}.

{\footnotesize\setlength{\tabcolsep}{4pt}\renewcommand{\arraystretch}{1.15}
\begin{table}[H]
\caption{Топ-5 лидеров по очкам и передачам (сезон 2025-26)}
\label{tab:sql-union-result}
\centering
\begin{tabular}{|l|l|r|}
\hline
\textbf{Категория} & \textbf{Игрок} & \textbf{Значение} \\
\hline
Scoring & Luka Dončić               & 32.97 \\
Scoring & Shai Gilgeous-Alexander   & 31.16 \\
Scoring & Anthony Edwards           & 28.80 \\
Scoring & Tyrese Maxey              & 28.29 \\
Scoring & Kawhi Leonard             & 27.89 \\
Assists & Nikola Jokić              & 10.72 \\
Assists & Cade Cunningham           &  9.63 \\
Assists & Josh Giddey               &  8.23 \\
Assists & Luka Dončić               &  8.15 \\
Assists & James Harden              &  8.14 \\
\hline
\end{tabular}
\end{table}}

Приведённые запросы охватывают простую фильтрацию, многотабличные
соединения, агрегацию с группировкой, скалярные подзапросы и составные
запросы с объединением результатов.

\section{Реализация логики предметной области в СУБД}

\subsection{Хранимые функции расчёта метрик}

Логика расчёта спортивных метрик реализована в виде хранимых
PL/pgSQL-функций (Приложение~В). Функция \texttt{calculate\_ts}
включает программную защиту от деления на ноль (Приложение~В,
листинг~\ref{lst:ts}).

Все расчётные функции помечены квалификатором \texttt{STABLE}, что позволяет
оптимизатору запросов кэшировать результаты вызовов в рамках одного
оператора. Агрегация сезонной статистики выполняется
хранимой процедурой \texttt{update\_season\_stats} с использованием
атомарной конструкции \texttt{INSERT ... ON CONFLICT DO UPDATE}.

\subsection{Триггеры логической целостности}

Для защиты данных от аномалий используется триггер
\texttt{trg\_check\_team\_score}. Квалификаторы \texttt{DEFERRABLE} и
\texttt{INITIALLY DEFERRED} задают отложенную проверку на этапе
\texttt{COMMIT}. Триггер проверяет бизнес-правило предметной области:
сумма очков всех игроков команды в конкретном матче должна равняться
итоговому командному счёту. Текст функции-триггера приведён
в приложении~В (листинг~\ref{lst:trigger}); команда создания~---
в листинге~\ref{lst:trigger-ddl}.

Триггер \texttt{trigger\_update\_season\_stats} с переходной таблицей
\texttt{inserted\_rows} обеспечивает автоматический пересчёт витрины
\texttt{player\_season\_stats} (листинг~\ref{lst:season-stats-trigger}).

\section{Реализация ролевой модели доступа}

Настройка прав доступа выполнена в соответствии с
принципом минимальных привилегий (Приложение~В, листинг~\ref{lst:grants}).
Роль публичного доступа (\texttt{reader})
лишена возможности прямого обращения к базовым таблицам; чтение
осуществляется исключительно через защищённые представления.

Корректность ограничений для роли \texttt{reader} проверена подключением к БД
от имени пользователя с назначенной ролью \texttt{reader}: запрос \texttt{SELECT * FROM game\_player\_stats} завершился
ошибкой \texttt{permission denied} на уровне PostgreSQL, тогда как
\texttt{SELECT} из представлений \texttt{v\_team\_stats} и \texttt{v\_top\_players},
перечисленных в приложении~В (листинг~\ref{lst:grants}), выполняется успешно.
Аналитические маршруты API, в том числе \texttt{GET /teams/standings},
открывают сессию под ролью \texttt{analyst} и для сводной командной
статистики обращаются к представлению \texttt{v\_team\_stats}.
У роли \texttt{analyst} есть \texttt{SELECT} к базовым таблицам (табл.~\ref{tab:acl-matrix}),
поэтому представления для неё служат удобными витринами, а не единственным
способом чтения данных.

\section{Реализация программного обеспечения и интерфейса}

\subsection{Серверное приложение}

Серверное приложение реализовано в соответствии со спецификацией REST API
конструкторского раздела (подраздел~\ref{subsec:api-spec}). Точкой входа
служит \texttt{app/main.py}, регистрирующий пять групп маршрутов и
CORS-middleware. Подключение к PostgreSQL организовано через три независимых
пула соединений~--- по одному на каждую роль СУБД. Привязка маршрута к
роли осуществляется через аргумент \texttt{Depends};
пример приведён в листинге~\ref{lst:di-role}:

\begin{lstlisting}[language=Python, label=lst:di-role,
  caption=Привязка маршрута к роли СУБД]
@router.get("/leaders")
async def get_leaders(
    season_id: int,
    db: AsyncSession = Depends(get_db_analyst),
    cache: CacheManager = Depends(get_cache),
):
\end{lstlisting}

Маршрут \texttt{/stats/leaders} открывает соединение от имени роли
\texttt{analyst}, имеющей право \texttt{SELECT} к агрегированной статистике
и \texttt{EXECUTE} к функциям расчёта метрик; прямое изменение базовых
таблиц по-прежнему запрещено. Жизненный цикл сессии управляется
генератором: при любом исходе обработчика сессия закрывается,
а незафиксированные изменения откатываются.

\subsection{Кэширование}

Алгоритм паттерна Cache-Aside (класс \texttt{CacheManager}): генерация ключа
на основе параметров, чтение из Redis; при промахе~--- выполнение SQL-запроса,
кэширование результата с установкой времени жизни (TTL) и возврат ответа.
Пример для маршрута \texttt{/stats/leaders} приведён в листинге~\ref{lst:cache-aside-pattern};
обобщённый вариант декоратора Cache-Aside~--- в Приложении~Г.

\begin{lstlisting}[language=Python, label=lst:cache-aside-pattern,
  caption=Cache-Aside в маршруте \texttt{/stats/leaders}]
cache_key = (
    f"leaders:{metric}:{season_id}:{team_id}:{position}:{min_games}:{limit}"
)
cached = await cache.get(cache_key)
if cached:
    return cached
result = await db.execute(query)
data = [...]
await cache.set(cache_key, data, ttl=900)
return data
\end{lstlisting}

Время жизни ключей дифференцировано по типу данных: данные завершённых
матчей кэшируются на один год, лидерборды и списки составов команд~---
на 15 минут,
результаты поиска~--- на 2 минуты. Инвалидация при вызове
\texttt{POST /admin/refresh/\{season\_id\}} выполняется по шаблонам ключей
затронутого сезона в соответствии с проектным решением
подраздела~\ref{subsec:cache}.

\subsection{Пользовательский интерфейс}

Интеграция библиотеки Recharts~\cite{recharts} обеспечивает рендеринг
диаграмм и графиков трендов.

Навигация включает аналитический дашборд с KPI-метриками лиги, таблицу
игроков с серверной фильтрацией, турнирные сводки конференций, глобальный
рейтинг лидеров и раздел расширенной статистики с диаграммами рассеяния
(рисунок~\ref{fig:ui-dashboard}). По выбору строки открывается модальное окно
профиля с сезонными метриками и игровым журналом (рисунок~\ref{fig:ui-modal}).
Модуль сравнения двух игроков реализован в виде модального окна
с радарным графиком, гистограммой базовой статистики и цветовым выделением
превосходящих метрик (рисунок~\ref{fig:ui-compare}).

\section{Тестирование и верификация аналитических данных}

Тестирование логики СУБД проводилось методом чёрного ящика. Для проверки
точности математических агрегаций выполнен регрессионный тест на реальных
данных из загруженной базы: эталонные значения получены независимым
расчётом по формулам раздела~1.2, после чего сопоставлены с результатами
хранимых функций PostgreSQL.

В качестве тестового объекта выбран профиль игрока Nikola Joki\'{c}
с идентификатором~2327 за регулярный сезон 2025-26 с идентификатором~5.
Прямой SQL-запрос суммы сырых фактов из таблицы
\texttt{game\_player\_stats} вернул следующие значения:
\begin{itemize}
    \item очки: 1799;
    \item броски с игры: 1132;
    \item штрафные попытки: 480.
\end{itemize}

Ручной расчёт истинного процента бросков по формуле~(\ref{eq:ts}):
\[
\mathrm{TS\%} = \frac{1799}{2 \cdot (1132 + 0{,}44 \cdot 480)}
  = \frac{1799}{2686{,}4} \approx 0{,}6697.
\]

Для верификации в СУБД выполнен вызов хранимой функции:
\begin{lstlisting}[language=SQL]
SELECT calculate_ts(2327, 5);
\end{lstlisting}

Результат совпал с ручным расчётом. Аналогичная верификация проведена для
eFG\% и USG\%; расхождений не выявлено.

Верификация PER выполнена по той же методике. Из таблицы
\texttt{game\_player\_stats} получены сезонные суммы для
\texttt{player\_id = 2327}: $\mathrm{PTS} = 1799$, $\mathrm{TRB} = 836$,
$\mathrm{AST} = 697$, $\mathrm{STL} = 92$, $\mathrm{BLK} = 53$,
$\mathrm{TOV} = 243$, $\mathrm{PF} = 173$, $\mathrm{FGM/FGA} = 644/1132$,
$\mathrm{FTM/FTA} = 399/480$, $\mathrm{MP} = 2264{,}8$.
Ненормализованный вклад по формуле uPER (раздел~1.2):
\[
\mathrm{uPER} =
(1799 + 836 + 697 + 92 + 53 - 243 - 173 - 488 - 35{,}64)
\cdot \frac{48}{2264{,}8} = 2537{,}36 \cdot \frac{48}{2264{,}8}
\approx 53{,}78.
\]
Лиговое среднее по формуле~(1.5) по 503~игрокам с $\mathrm{MP} > 100$:
$\mathrm{lg\_aPER} = 21{,}47$. Итоговый рейтинг по формуле~(1.6):
\[
\mathrm{PER} = 53{,}78 \cdot \frac{15}{21{,}47} \approx 37{,}57.
\]
Вызов \texttt{SELECT calculate\_per(2327, 5)} вернул $37{,}57$~---
расхождений не выявлено.

Тестирование граничных случаев подтвердило корректный возврат значения
\texttt{NULL} при отсутствии попыток бросков и нулевом времени на площадке.
Работа триггера \newline\texttt{trg\_check\_team\_score} проверена попыткой вставки
в одной транзакции строк с намеренно неверной суммой очков: при
\texttt{COMMIT} триггер поднял исключение
\texttt{Player points sum does not match team score}, транзакция была
отклонена~--- поведение соответствует ожидаемому
(таблица~\ref{tab:test-cases}).

{\footnotesize\setlength{\tabcolsep}{3pt}\renewcommand{\arraystretch}{1.2}
\begin{table}[H]
\caption{Контрольные тест-кейсы функций расчёта метрик}
\label{tab:test-cases}
\centering
\begin{tabular}{|c|p{0.24\textwidth}|p{0.27\textwidth}|p{0.17\textwidth}|p{0.14\textwidth}|}
\hline
\textbf{№} & \textbf{Сценарий} & \textbf{Входные данные}
  & \textbf{Ожидание} & \textbf{Результат} \\
\hline
1 & Нормальный профиль & \texttt{calculate\_ts(2327, 5)}
  & $\approx 0{,}6697$ & совпало \\
\hline
2 & Нет попыток броска & \texttt{calculate\_ts(0, 0)}
  & NULL & NULL \\
\hline
3 & Нулевые FGA (eFG\%) & \texttt{calculate\_efg(0, 0)}
  & NULL & NULL \\
\hline
4 & Пересчёт витрины & \texttt{CALL update\_season\_stats(5)}
  & обновлены строки \texttt{pss} & выполнено \\
\hline
5 & Верификация PER & \texttt{calculate\_per(2327, 5)}
  & $\approx 37{,}57$ (формулы 1.4--1.6) & 37{,}57, совпало \\
\hline
6 & Триггер целостности счёта & INSERT с суммой очков,
  не совпадающей с \texttt{home\_score}
  & EXCEPTION при попытке фиксации транзакции (COMMIT) & исключение при COMMIT,
  транзакция отклонена \\
\hline
\end{tabular}
\end{table}
}

Для формализации тестового покрытия каждой функции на стороне СУБД
выделены классы эквивалентности~--- непересекающиеся разделы входных
данных, для которых ожидается одинаковое поведение функции. Разбиение
приведено в таблице~\ref{tab:eq-classes}.

\begin{table}[H]
\caption{Классы эквивалентности функций на стороне СУБД}
\label{tab:eq-classes}
\centering
\begin{tabular}{|>{\raggedright\arraybackslash}p{0.21\textwidth}|c|>{\raggedright\arraybackslash}p{0.28\textwidth}|>{\raggedright\arraybackslash}p{0.22\textwidth}|>{\raggedright\arraybackslash}p{0.14\textwidth}|}
\hline
\textbf{Функция} & \textbf{КЭ} & \textbf{Условие раздела} & \textbf{Представитель} & \textbf{Ожидание} \\
\hline

\texttt{calculate\_efg}
 & КЭ1 & FGA $> 0$, FG3M $\geq 0$ & \texttt{calc\_efg(2327, 5)} & $\in (0; 1]$ \\
 & КЭ2 & FGA $= 0$ (нет бросков) & \texttt{calc\_efg(0, 0)} & NULL \\
 & КЭ3 & team\_id IS NULL (агрегация по всем командам) & \texttt{calc\_efg(2327, 5, NULL)} & $=$ КЭ1 \\
\hline
\texttt{calculate\_usg}
 & КЭ1 & MP $\geq 1$, denom $> 0$ & \texttt{calc\_usg(2327, 5)} & $\in (0; 1]$ \\
 & КЭ2 & MP $< 1$ (игрок не выходил) & несуществующий player\_id & NULL \\
 & КЭ3 & командный знаменатель $= 0$ & аномальные тестовые данные & NULL \\
 & КЭ4 & team\_id не определён (нет матчей) & несуществующий player\_id & NULL \\
\hline
\texttt{calculate\_ts}
 & КЭ1 & FGA $> 0$ или FTA $> 0$ & \texttt{calc\_ts(2327, 5)} & $\approx 0{,}6697$ \\
 & КЭ2 & FGA $= 0$ и FTA $= 0$ & \texttt{calc\_ts(0, 0)} & NULL \\
 & КЭ3 & знаменатель $= 0$ при FGA $= 0$, FTA $> 0$ (редкий) & аномальные тестовые данные & NULL \\
\hline
\texttt{calculate\_per}
 & КЭ1 & MP $\geq 50$, lg\_aPER $> 0$ & \texttt{calc\_per(2327, 5)} & $\approx 37{,}57$ \\
 & КЭ2 & MP $< 50$ (дисквалифицирован) & игрок с малым числом игр & NULL \\
 & КЭ3 & lg\_aPER $= 0$ или NULL (пустой сезон) & фиктивный season\_id без данных & нормировка на $15{,}0$ \\
\hline
\texttt{calculate\_bpm}
 & КЭ1 & MP $\geq 50$ & \texttt{calc\_bpm(2327, 5)} & числовое значение \\
 & КЭ2 & MP $< 50$ & игрок с малым числом игр & NULL \\
\hline
\texttt{trg\_validate\_} \texttt{team\_score\_row}
 & КЭ1 & статус $\neq$ \texttt{Finished} или счёт NULL & INSERT в незавершённый матч & принято без проверки \\
 & КЭ2 & статус \texttt{Finished}, $\sum\mathrm{points} =$ счёт & корректный INSERT & принято \\
 & КЭ3 & статус \texttt{Finished}, $\sum\mathrm{points} \neq$ счёт & INSERT с неверной суммой очков & EXCEPTION \\
\hline
\texttt{trigger\_update\_} \texttt{season\_stats}
 & КЭ1 & season\_id найден & штатный INSERT в \texttt{game\_player\_stats} & вызов \texttt{update\_} \texttt{season\_} \texttt{stats} \\
 & КЭ2 & season\_id IS NULL (пустая таблица) & таблица без строк & возврат NULL без вызова \\
\hline
\end{tabular}
\end{table}

Тестирование выполнено по всем выделенным классам эквивалентности.
Для каждого класса фактическое поведение функции совпало с ожидаемым:
корректные входные данные дают значения в допустимом диапазоне, граничные
и вырожденные случаи возвращают NULL либо предусмотренную нормировку,
а нарушение бизнес-правила счёта вызывает исключение на этапе фиксации
транзакции. Все тесты пройдены успешно, расхождений не выявлено.

\section*{Вывод}

В технологическом разделе обоснован стек технологий (PostgreSQL~15,
Redis~7, FastAPI, React~18) и реализована структура базы данных средствами
DDL: таблицы, ограничения целостности и индексы.
Разработаны PL/pgSQL-функции и триггеры, обеспечивающие расчёт метрик и
консистентность данных. С помощью написанного ETL-скрипта БД наполнена
реальной статистикой NBA (более 150~тыс. фактов). Успешно реализованы
механизмы кэширования (Cache-Aside) и ролевого доступа. Регрессионное
тестирование на реальных данных подтвердило математическую корректность
работы базы данных и хранимых процедур.